\documentclass[12pt]{article}

\usepackage[a4paper,margin=1in]{geometry}
\usepackage[T1]{fontenc}
\usepackage[utf8]{inputenc}
\usepackage{lmodern}
\usepackage{amsmath,amssymb}
\usepackage{booktabs}
\usepackage{longtable}
\usepackage{array}
\usepackage{tabularx}
\usepackage{enumitem}
\usepackage{hyperref}
\usepackage{setspace}
\usepackage{ragged2e}

\setstretch{1.08}
\setlength{\parskip}{0.5em}
\setlength{\parindent}{0pt}

\hypersetup{
	colorlinks=true,
	linkcolor=black,
	citecolor=black,
	urlcolor=blue,
	pdftitle={Proof Status Register for the General Theory of Cognitive Structuring},
	pdfauthor={Kostiantyn Osmolovskyi}
}

\title{Proof Status Register\\
	\small for the General Theory of Cognitive Structuring}

\author{
	Kostiantyn Osmolovskyi\\
	\small Independent Researcher, Ukraine\\
	\small ORCID: 0009-0006-3144-7237\\
	\small \texttt{constantinosmol@gmail.com}
}

\date{}

% ---------------------------------------------------------------------------
% Part of a series: General Theory of Cognitive Structuring (GTCS)
% Autor: Kostiantyn Osmolovskyi (ORCID: 0009-0006-3144-7237)
% License: Creative Commons Attribution 4.0 International (CC BY 4.0)
% Repository: https://github.com/k-osmolovskyi/k-osmolovskyi.github.io
% DOI: https://doi.org/10.5281/zenodo.19705732
% ---------------------------------------------------------------------------

\begin{document}
	\maketitle

\begin{center}
	\small
	Verification Report No. 2026-3\\
	Version 1.0
\end{center}

\vspace{0.5em}	
	
	\section{Purpose of this document}
	
	This document provides a compact proof-status register for the \emph{General Theory of
		Cognitive Structuring}. Its purpose is to identify, in one place, the current proof status of
	the main results of the series and to distinguish clearly between:
	\begin{itemize}[leftmargin=2em]
		\item results with full proof in source,
		\item results supported by proof sketch,
		\item canonically normalized results,
		\item architectural or definitional inclusions,
		\item and results pending later consolidation.
	\end{itemize}
	
	The register is intended to serve as an intermediate canonical layer between:
	\begin{itemize}[leftmargin=2em]
		\item the \emph{Collected Results and Proof Status Note},
		\item the \emph{Collected Propositions and Theorems},
		\item and a future full \emph{Proof Compendium}.
	\end{itemize}
	
	It is not itself a proof compendium and does not reproduce proofs in full.
	
	\section{Scope and limits}
	
	This register is limited to proof status and proof-readiness.
	
	It does \emph{not} attempt to:
	\begin{itemize}[leftmargin=2em]
		\item restate the full theory,
		\item reproduce the detailed proofs,
		\item replace the source papers,
		\item function as a complete theorem numbering system,
		\item resolve every notational or editorial variation in the corpus.
	\end{itemize}
	
	Its narrower role is:
	\begin{itemize}[leftmargin=2em]
		\item to provide a stable overview of which results are already proof-complete,
		\item to identify which results remain sketch-level,
		\item to identify which results are canonical but not theorem-grade,
		\item to mark which results are priority targets for later proof extraction.
	\end{itemize}
	
	\section{Proof status categories}
	
	\begin{longtable}{>{\RaggedRight\arraybackslash}p{0.10\textwidth}
			>{\RaggedRight\arraybackslash}p{0.75\textwidth}}
		\toprule
		\textbf{Code} & \textbf{Meaning} \\
		\midrule
		\endfirsthead
		\toprule
		\textbf{Code} & \textbf{Meaning} \\
		\midrule
		\endhead
		
		D & Definitional item. No proof is required because the item functions as a definition, operator, or construction. \\
		
		A & Architectural claim. The item is central to the architectural interpretation of the framework but is not yet treated as a theorem-grade result. \\
		
		P & Full proof in source. The result is stated and proved in the source paper in a form sufficient for direct canonical reference. \\
		
		S & Proof sketch in source. The result is supported in the source paper but remains compressed, local, or sketch-level in its current canonical form. \\
		
		C & Canonically normalized. The result is collected and normalized across one or more papers, but this register does not itself upgrade it to a full proof if such proof is not already available in source. \\
		
		PC & Proof-complete after normalization. The source proof is sufficiently complete and the canonical normalized statement may be used as proof-ready without major further consolidation. \\
		
		PP & Proof-pending consolidation. The result is important enough for canonical inclusion but still requires later extraction, normalization, or auxiliary lemma support before proof-compendium use. \\
		
		O & Open or reserved for future proof-compendium treatment. \\
		
		\bottomrule
	\end{longtable}
	
	\section{How to read this register}
	
	Each item in the register includes:
	\begin{itemize}[leftmargin=2em]
		\item \textbf{Result ID}
		\item \textbf{Title}
		\item \textbf{Type}
		\item \textbf{Source}
		\item \textbf{Current status}
		\item \textbf{Proof basis}
		\item \textbf{Canonical use}
		\item \textbf{Next step}
	\end{itemize}
	
	The register should be read together with:
	\begin{itemize}[leftmargin=2em]
		\item the \emph{Collected Propositions and Theorems},
		\item the \emph{Collected Results and Proof Status Note},
		\item and the \emph{Technical Appendix}.
	\end{itemize}
	
	\section{Master proof status register}
	
	\begin{longtable}{>{\RaggedRight\arraybackslash}p{0.07\textwidth}
			>{\RaggedRight\arraybackslash}p{0.24\textwidth}
			>{\RaggedRight\arraybackslash}p{0.12\textwidth}
			>{\RaggedRight\arraybackslash}p{0.14\textwidth}
			>{\RaggedRight\arraybackslash}p{0.08\textwidth}
			>{\RaggedRight\arraybackslash}p{0.16\textwidth}
			>{\RaggedRight\arraybackslash}p{0.10\textwidth}} 
		\toprule
		\textbf{Result ID} & \textbf{Title} & \textbf{Type} & \textbf{Source} & \textbf{Status} & \textbf{Proof basis} & \textbf{Next step} \\
		\midrule
		\endfirsthead
		\toprule
		\textbf{Result ID} & \textbf{Title} & \textbf{Type} & \textbf{Source} & \textbf{Status} & \textbf{Proof basis} & \textbf{Next step} \\
		\midrule
		\endhead
		
		G1.01 & Coherence as distance to the stability region & Definition-formal proposition & coherence paper & PC & full source proof and stable normalization & direct reuse \\
		
		G1.02 & Coherence is distinct from structural tension & Proposition & coherence paper & C & architectural distinction normalized from source & optional proof note \\
		
		R2.01 & Structural admissibility depends on the joint regulatory state & Proposition & structural admissibility & PC & full source proof and stable normalization & direct reuse \\
		
		R2.02 & Overload memory is not reducible to instantaneous tension & Proposition & overload formation & PC & full source proof & direct reuse \\
		
		R2.03 & Regulatory state is trajectory-dependent & Proposition & trajectory regulation & PC & full source proof & direct reuse \\
		
		R2.04 & Overload deforms the admissibility boundary and introduces hysteresis & Proposition & trajectory regulation & S/PP & source argument with explicit assumption and compressed theorem form & later normalized proof extraction \\
		
		R2.05 & The regulatory phase space is minimally two-dimensional & Proposition & trajectory regulation + earlier regulatory papers & C & canonical integration across papers & optional formal extraction \\
		
		R3.01 & Invariant-induced stability geometry & Proposition /Theorem & invariants paper & PC & full source proof & direct reuse \\
		
		R3.02 & Global conflict is structured by invariant interaction architecture & Proposition & invariants paper & PC & full source proof & direct reuse \\
		
		R3.03 & Compression is not equivalent to simplification & Proposition & compression paper & PC & full source proof & direct reuse \\
		
		R3.04 & Compression generates a new invariant & Proposition & compression paper & PC & full source proof & direct reuse \\
		
		R3.05 & Compression can preserve or transform admissible geometry & Proposition & compression paper & S/PP & source support clear but theorem normalization may be sharpened & later normalized extraction \\
		
		R3.06 & Level transition is a non-equivalent change in admissible geometry & Proposition & structural admissibility + compression & C/PP & distributed proof basis across papers & later collected proof \\
		
		R3.07 & Compression improves effective long-run regulatory stability & Proposition & compression paper & S/PP & source support tied to long-run criterion & later assumption-explicit proof \\
		
		R4.01 & Coherence representation is an internal regulatory variable rather than a reconstruction of geometric coherence & Proposition & emergence paper & PC & full source proof and stable role separation & direct reuse \\
		
		R4.02 & Coherence representation preserves regulatory order without requiring metric preservation & Proposition & emergence + multi-system regulation & C/PP & distributed source support & later collected proof \\
		
		R4.03 & Pre-symbolic admissibility is a distinct prior filtering layer & Proposition-like definitional result & pre-symbolic paper & D/C & definitional separation of layers & no proof required \\
		
		R4.04 & Pre-symbolic admissibility determines the admissible discrepancy domain & Proposition & pre-symbolic paper & S/PP & source support present but theorem form compressed & later normalized proof \\
		
		R4.05 & Restricted accessibility of coherence & Proposition & restricted accessibility paper & PC & full source proof & direct reuse \\
		
		R4.06 & Non-injectivity of regulatory accessibility & Proposition & restricted accessibility paper & S/PP & source proof basis present but canonical theorem extraction desirable & later normalized proof \\
		
		R4.07 & Restricted accessibility is structural, not merely representationally approximate & Proposition & restricted accessibility paper & C & canonical consequence from source argument & optional proof note \\
		
		R4.08 & Repeated admissible propagation induces propagation bias & Proposition & pre-symbolic paper & S/PP & source argument present at sketch level & later explicit proof extraction \\
		
		R4.09 & Repeated admissible propagation may stabilize proto-structural channels & Proposition & pre-symbolic paper & O/PP & important but still pending stronger canonical proof form & future compendium priority \\
		
		R5.01 & Identity is derived as a regulatory attractor rather than introduced as a primitive self-representational object & Proposition & identity paper & C/PP & canonical architectural normalization from source & later formal sharpening \\
		
		R5.02 & Low-overload concentration regions may function as identity-cores & Proposition /Theorem & identity paper & S/PP & source theorem family tied to long-run assumptions & later normalized theorem extraction \\
		
		R5.03 & Identity persistence is regulatory rather than merely representational & Proposition & identity paper & C & canonical architectural result & optional proof note \\
		
		R5.04 & Inter-system comparability does not require shared metric coherence or shared configuration space & Proposition & multi-system regulation & C/PP & source argument strong but best collected canonically & later normalization \\
		
		R5.05 & Minimal transferable structure between systems is ordinal rather than metric & Proposition & multi-system regulation & PC & full source proof & direct reuse \\
		
		R5.06 & Multi-system co-regulation operates through aligned regulatory variables rather than through shared global structure & Proposition & multi-system regulation & C/PP & canonical normalization from source support & later explicit proof extraction \\
		
		R5.07 & Inter-system conflict is incompatibility of admissible discrepancy domains sufficient for coordinated regulation & Proposition /Definitional result & inter-system conflict & S/PP & source support present but mixed definitional /propositional status & later theorem-grade clarification \\
		
		R5.08 & Inter-system conflict may exist without direct representation in one or both systems & Proposition & inter-system conflict & C/PP & canonical extension of restricted-accessibility logic & later explicit proof extraction \\
		
		R5.09 & Alignment and conflict are distinct extension problems & Proposition & multi-system regulation + inter-system conflict & C & synthetic canonical clarification across two papers & no immediate proof need \\
		
		R6.01 & The full framework is organized by layered admissibility & Synthetic proposition & general theory & C & canonical synthetic convergence statement & synthetic use only \\
		
		R6.02 & Structural inconsistency, discrepancy accessibility, representation, and structural updating are distinct layers & Synthetic proposition & general theory & C & canonical anti-collapse statement & synthetic use only \\
		
		R6.03 & Representation is defined over a restricted admissible domain and does not reconstruct full structural deviation & Synthetic proposition & general theory & C & canonical integration of access and representation branches & synthetic use only \\
		
		R6.04 & Overload introduces historical depth into admissibility across the framework & Synthetic proposition & general theory + overload / trajectory papers & C & canonical integration of regulatory branch & synthetic use only \\
		
		R6.05 & Invariant structure is the architectural source of admissible geometry and compression is the mechanism of its non-equivalent transformation & Synthetic proposition & general theory & C & canonical integration of invariant and compression branches & synthetic use only \\
		
		R6.06 & Identity, representation, and inter-system regulation are downstream consequences of a common layered architecture & Synthetic proposition & general theory & C & canonical cross-branch convergence statement & synthetic use only \\
		
		R6.07 & The General Theory functions as a synthetic convergence framework rather than as the hidden source of earlier branches & Meta-structural proposition & verification package + general theory & C & meta-structural canonical statement & verification use only \\
		
		\bottomrule
	\end{longtable}
	
	\section{Results ready for direct canonical reuse}
	
	The following results are already proof-complete enough to be used directly in a later canonical
	proof layer with only minimal normalization:
	
	\begin{itemize}[leftmargin=2em]
		\item G1.01 --- Coherence as distance to the stability region
		\item R2.01 --- Structural admissibility depends on the joint regulatory state
		\item R2.02 --- Overload memory is not reducible to instantaneous tension
		\item R2.03 --- Regulatory state is trajectory-dependent
		\item R3.01 --- Invariant-induced stability geometry
		\item R3.02 --- Global conflict is structured by invariant interaction architecture
		\item R3.03 --- Compression is not equivalent to simplification
		\item R3.04 --- Compression generates a new invariant
		\item R4.01 --- Coherence representation is an internal regulatory variable rather than a reconstruction of geometric coherence
		\item R4.05 --- Restricted accessibility of coherence
		\item R5.05 --- Minimal transferable structure between systems is ordinal rather than metric
	\end{itemize}
	
	\section{Results requiring later normalized extraction}
	
	The following results are central but should be normalized more explicitly before being treated
	as fully proof-ready in a future compendium:
	
	\begin{itemize}[leftmargin=2em]
		\item R2.04 --- Overload deforms the admissibility boundary and introduces hysteresis
		\item R3.05 --- Compression can preserve or transform admissible geometry
		\item R3.06 --- Level transition is a non-equivalent change in admissible geometry
		\item R3.07 --- Compression improves effective long-run regulatory stability
		\item R4.02 --- Coherence representation preserves regulatory order without requiring metric preservation
		\item R4.04 --- Pre-symbolic admissibility determines the admissible discrepancy domain
		\item R4.06 --- Non-injectivity of regulatory accessibility
		\item R4.08 --- Repeated admissible propagation induces propagation bias
		\item R5.01 --- Identity is derived as a regulatory attractor rather than introduced as a primitive self-representational object
		\item R5.02 --- Low-overload concentration regions may function as identity-cores
		\item R5.04 --- Inter-system comparability does not require shared metric coherence or shared configuration space
		\item R5.06 --- Multi-system co-regulation operates through aligned regulatory variables rather than through shared global structure
		\item R5.07 --- Inter-system conflict is incompatibility of admissible discrepancy domains sufficient for coordinated regulation
		\item R5.08 --- Inter-system conflict may exist without direct representation in one or both systems
	\end{itemize}
	
	\section{Results pending future proof-compendium treatment}
	
	The following results are important and already canonically recognized, but they remain priority
	targets for later proof-compendium treatment:
	
	\begin{itemize}[leftmargin=2em]
		\item R4.09 --- Repeated admissible propagation may stabilize proto-structural channels
		\item R5.02 --- Low-overload concentration regions may function as identity-cores
		\item R5.07 --- Inter-system conflict is incompatibility of admissible discrepancy domains sufficient for coordinated regulation
	\end{itemize}
	
	\section{Canonical synthetic results}
	
	The following items are canonically included in the framework but should remain explicitly
	synthetic or meta-structural rather than theorem-foundational in a later proof architecture:
	
	\begin{itemize}[leftmargin=2em]
		\item R6.01 --- The full framework is organized by layered admissibility
		\item R6.02 --- Structural inconsistency, discrepancy accessibility, representation, and structural updating are distinct layers
		\item R6.03 --- Representation is defined over a restricted admissible domain and does not reconstruct full structural deviation
		\item R6.04 --- Overload introduces historical depth into admissibility across the framework
		\item R6.05 --- Invariant structure is the architectural source of admissible geometry and compression is the mechanism of its non-equivalent transformation
		\item R6.06 --- Identity, representation, and inter-system regulation are downstream consequences of a common layered architecture
		\item R6.07 --- The General Theory functions as a synthetic convergence framework rather than as the hidden source of earlier branches
	\end{itemize}
	
	\section{Priority list for a future proof compendium}
	
	A future canonical proof compendium would most naturally begin with the following sequence:
	
	\subsection*{Stage 1: Core geometric and regulatory backbone}
	\begin{itemize}[leftmargin=2em]
		\item G1.01
		\item R2.01
		\item R2.02
		\item R2.03
		\item R2.04
	\end{itemize}
	
	\subsection*{Stage 2: Invariant and compression branch}
	\begin{itemize}[leftmargin=2em]
		\item R3.01
		\item R3.02
		\item R3.03
		\item R3.04
		\item R3.05
		\item R3.06
	\end{itemize}
	
	\subsection*{Stage 3: Representation and accessibility branch}
	\begin{itemize}[leftmargin=2em]
		\item R4.01
		\item R4.02
		\item R4.04
		\item R4.05
		\item R4.06
	\end{itemize}
	
	\subsection*{Stage 4: Identity and multi-system branch}
	\begin{itemize}[leftmargin=2em]
		\item R5.01
		\item R5.02
		\item R5.05
		\item R5.07
		\item R5.08
	\end{itemize}
	
	\section{Compact proof-status summary table}
	
	\begin{longtable}{>{\RaggedRight\arraybackslash}p{0.10\textwidth}
			>{\RaggedRight\arraybackslash}p{0.44\textwidth}
			>{\RaggedRight\arraybackslash}p{0.12\textwidth}
			>{\RaggedRight\arraybackslash}p{0.18\textwidth}}
		\toprule
		\textbf{ID} & \textbf{Title} & \textbf{Status} & \textbf{Immediate use} \\
		\midrule
		\endfirsthead
		\toprule
		\textbf{ID} & \textbf{Title} & \textbf{Status} & \textbf{Immediate use} \\
		\midrule
		\endhead
		
		G1.01 & Coherence as distance to the stability region & PC & proof-ready \\
		
		R2.01 & Structural admissibility depends on the joint regulatory state & PC & proof-ready \\
		
		R2.02 & Overload memory is not reducible to instantaneous tension & PC & proof-ready \\
		
		R2.03 & Regulatory state is trajectory-dependent & PC & proof-ready \\
		
		R2.04 & Overload deforms the admissibility boundary and introduces hysteresis & S/PP & normalize later \\
		
		R3.01 & Invariant-induced stability geometry & PC & proof-ready \\
		
		R3.03 & Compression is not equivalent to simplification & PC & proof-ready \\
		
		R3.04 & Compression generates a new invariant & PC & proof-ready \\
		
		R3.05 & Compression can preserve or transform admissible geometry & S/PP & normalize later \\
		
		R4.01 & Coherence representation is an internal regulatory variable rather than a reconstruction of geometric coherence & PC & proof-ready \\
		
		R4.05 & Restricted accessibility of coherence & PC & proof-ready \\
		
		R4.06 & Non-injectivity of regulatory accessibility & S/PP & normalize later \\
		
		R5.02 & Low-overload concentration regions may function as identity-cores & S/PP & compendium priority \\
		
		R5.05 & Minimal transferable structure between systems is ordinal rather than metric & PC & proof-ready \\
		
		R5.07 & Inter-system conflict is incompatibility of admissible discrepancy domains sufficient for coordinated regulation & S/PP & compendium priority \\
		
		R6.01 & The full framework is organized by layered admissibility & C & synthetic canonical use \\
		
		\bottomrule
	\end{longtable}
	
	\section{Closing note}
	
	The purpose of this register is not to turn the present framework prematurely into a monolithic proof architecture. Its role is to fix the current canonical author-side assessment of which results are already proof-ready, which remain sketch-level, which are synthetic, and which are the highest priorities for future proof-compendium treatment.
	
\end{document}